\documentclass[11pt]{article}
\usepackage[T1]{fontenc}
\usepackage[utf8]{inputenc}
\usepackage{amsmath,amssymb,amsthm,booktabs}
\usepackage[margin=1in]{geometry}
\usepackage{hyperref,xurl}
\newtheorem{theorem}{Theorem}[section]
\newtheorem{lemma}[theorem]{Lemma}
\newtheorem{proposition}[theorem]{Proposition}
\newtheorem{remark}[theorem]{Remark}
\newcommand{\HH}{\mathbb H}
\newcommand{\Vol}{\operatorname{Vol}}
\newcommand{\ip}[2]{\langle #1,#2\rangle}
\title{The second distinct volume of entirely ideal\linebreak right-angled hyperbolic four-polytopes}
\author{Coxeter Polytopes Frontier Project}
\date{Proof translation, 18 October 2026}
\begin{document}
\maketitle
\begin{abstract}
Among finite-volume entirely ideal right-angled hyperbolic four-polytopes,
the smallest volume is $4\pi^2/3$, attained uniquely by the ideal 24-cell,
by Kolpakov. We prove that the second distinct volume is $8\pi^2/3$.
The new lower bound uses facet defects, the geometry of opposite faces
of a cusp box, and the published classification of the nine-vertex
ideal right-angled three-polyhedron. Reflection of the 24-cell across
one facet attains the bound. An exact rational certificate checks this
double and its complete face and cusp data. The universal geometric
argument is a written proof, not a finite four-polytope census.
This manuscript translates the internally checked argument; external
refereeing and an end-to-end formalization are not claimed.
Classification of all polytopes attaining the second value is not a
theorem of this manuscript.
\end{abstract}

\section{Statement and inputs}
Throughout, a polytope is finite-sided, convex, four-dimensional, and of
finite hyperbolic volume. ``Entirely ideal'' means that every vertex is
on the sphere at infinity. ``Right-angled'' means that every dihedral
angle is $\pi/2$. Let $C$ denote the regular ideal right-angled 24-cell.

\begin{theorem}\label{thm:gap}
If $P$ is an entirely ideal right-angled hyperbolic four-polytope and
$P$ is not isometric to $C$, then
\[
 \Vol(P)\ge \frac{8\pi^2}{3}.
\]
Equality is attained by the union of $C$ and its reflection across an
octahedral facet. Consequently the second distinct volume is
$8\pi^2/3$.
\end{theorem}

The distinction between a second volume \emph{value} and the
classification of its equality cases is essential. Theorem~\ref{thm:gap}
does not assert uniqueness at $8\pi^2/3$.

We use the following published inputs, with the indicated precise scope.
Kolpakov's Lemmas~1--2 give the four-dimensional incidence and volume
formulas, and Theorem~3 identifies $C$ as the unique minimum
\cite{Kolpakov}. Proposition~5 states that an ideal right-angled
three-polyhedron with a $k$-gonal face has at least $2k$ vertices, with
equality exactly for the $k$-antiprism. Section~4.3, Case~D, and Figure~6
give the complete admissible small octahedrite list and, in particular,
the unique nine-vertex case. We use only that nine-vertex classification,
not a classification of arbitrary four-polytopes. The standard
right-angled rigidity of these three-dimensional polyhedra identifies
the classified combinatorial types with their geometric models, as in
Kolpakov's use of Andreev's theorem.

The new opposite-cusp injection, cusp-shape calculation, minimum-defect
argument and their combination are internal arguments below. They are
not treated as published assumptions. The exact construction certificate
has a separate scope from the universal proof.

\section{Discrete volume and facet defect}
Write $v,e,r,f$ for the numbers of vertices, edges, ridges and facets of
$P$. An ideal vertex link is a Euclidean rectangular three-box. Counting
its vertices and edges and using Euler's identity gives
\[
 e=4v,\qquad r=3v+f,\qquad
 \sum_{R\in\Omega_2(P)} k(R)=12v,
\]
where $k(R)$ is the number of vertices of a ridge $R$. Define
\[
 W(P)=\sum_R(k(R)-3).
\]
It follows that
\begin{equation}\label{eq:discrete}
 W(P)=3(v-f)\in3\mathbb Z_{\ge0},\qquad
 \Vol(P)=\frac{\pi^2}{3}(v-f+4)
 =\frac{4\pi^2}{3}+\frac{\pi^2}{9}W(P).
\end{equation}
The latter identity is Kolpakov's volume formula; equivalently,
Gauss--Bonnet applied to the reflection orbifold uses
$\chi=1-f/2+r/4-e/8=(v-f+4)/4$.

Every facet $F$ is an ideal right-angled three-polyhedron. Its ideal
vertices have valency four, so, writing $v_F$ for their number, it has
$2v_F$ edges and $v_F+2$ polygonal faces. Its defect is
\begin{equation}\label{eq:facet}
 d(F):=v_F-6=\sum_{R\subset F}(k(R)-3),\qquad
 \sum_F d(F)=2W(P).
\end{equation}
Defect zero characterizes the regular ideal octahedron. A positive
defect facet has a nontriangular face of size $k\ge4$; the antiprism
bound $v_F\ge2k$ implies $d(F)\ge2$. In particular defect one does not
occur.

\section{Opposite cusp facets}
For a facet $F$, let $q(F)$ be its number of nontriangular faces and
$s(F)$ the number of its cusps whose intrinsic horospherical rectangle
is nonsquare. The latter condition does not depend on horosphere scale.

\begin{lemma}[Opposite-cusp injection]\label{lem:injection}
Suppose the smallest positive facet defect of $P$ is $m$, and $F$ has
positive defect. There are $q(F)+s(F)$ pairwise distinct positive-defect
facets other than $F$: its neighbors across nontriangular ridges, and
the facets opposite $F$ at its nonsquare cusps. Therefore
\begin{equation}\label{eq:budget}
 2W(P)\ge d(F)+m\bigl(q(F)+s(F)\bigr).
\end{equation}
In particular, when $d(F)=m$,
$2W(P)\ge m(1+q(F)+s(F))$.
\end{lemma}
\begin{proof}
A neighbor across a nontriangle contains that same nontriangle and is
not octahedral. Distinct ridges of $F$ have distinct neighboring facets,
since the intersection of two facets is a face.

At a cusp $x$ of $F$, the face of the rectangular three-box opposite
the $F$-face belongs to a facet $F_x$. The two rectangles are congruent.
If the rectangle of $F$ is nonsquare, so is that of $F_x$, which cannot
be the regular ideal octahedron.

The supporting hyperplanes of $F$ and $F_x$ are asymptotic and meet
at exactly the ideal point $x$. Hence the same $F_x$ cannot occur at
a second cusp of $F$, nor can it meet $F$ across any ridge. Thus all
facets in the two lists are distinct and differ from $F$. Each has
defect at least $m$. Summing their defects in \eqref{eq:facet} proves
the assertion.
\end{proof}

If $P\not\cong C$, a positive-defect facet exists: otherwise $W=0$
in \eqref{eq:facet}, and the volume formula and Kolpakov's unique
minimum identify $P$ with $C$.

\section{The two cusp-shape calculations}
\begin{lemma}\label{lem:antiprism}
The cusp rectangles of the ideal right-angled $k$-antiprism have
aspect ratio $2\cos(\pi/k)$. Thus for $k\ge4$ all $2k$ cusps
are nonsquare.
\end{lemma}
\begin{proof}
Set $\theta=\pi/k$, $c=\cos\theta$, and take positive square roots in
\[
 h^2=\frac{1-c}{1+c},\quad r^2=\frac{2c}{1+c},\quad
 b^2=\frac{2}{c(1+c)},\quad \mu^2=\frac{2(1-c)}{c}.
\]
Use Lorentz form $(-1,1,1,1)$. The two vertex rings have heights
$\pm h$, radius $r$, and angular offset $\theta$. At
$u=(1,r\cos\theta,r\sin\theta,h)$, the four incident normals are
\[
 N=(h,0,0,1),\quad M_0=(1,b,0,h),\quad
 M_2=(1,b\cos2\theta,b\sin2\theta,h),\quad
 L=(1,b\cos\theta,b\sin\theta,-h).
\]
Their squares are $r^2,\mu^2,\mu^2,\mu^2$. Adjacent products vanish;
the opposite pairs are parallel. The identities $bc=r$ and
$\mu=2h/r$ give
\[
 \frac{M_0+M_2}{\mu}=\frac2\mu u,\qquad
 \frac Nr+\frac L\mu=\left(\frac hr+\frac1\mu\right)u.
\]
For two opposite unit normals with sum $\kappa u$, their separation
on a fixed horosphere is proportional to $\kappa$. To see this,
send $u$ to infinity in the upper-half-space model; the mirrors become
parallel vertical planes with unit horizontal normals. The two
separations therefore have ratio
\[
 \frac{2/\mu}{h/r+1/\mu}
 =\frac{2r^2}{1+h^2}=2c.
\]
The antiprism symmetries are vertex-transitive. Finally $2c>1$ for
$k\ge4$ (and equals one for $k=3$).
\end{proof}

\begin{lemma}\label{lem:nine}
An ideal right-angled three-polyhedron with nine vertices has
$d=3$, $q=3$ and $s=3$.
\end{lemma}
\begin{proof}
The antiprism bound excludes pentagons and larger faces. The complete
published nine-vertex octahedrite list consists of one admissible type:
two ideal regular octahedra glued across a triangular face
\cite[Section~4.3, Case~D2 and Figure~6]{Kolpakov}. Rigidity makes
this identification geometric. The three gluing cusps have rectangles
formed by joining two squares, hence aspect two. The other six cusps
remain square. Its three nontriangular faces are quadrilaterals.
\end{proof}

\section{Universal exclusion of intermediate volumes}
\begin{lemma}\label{lem:two}
A positive-defect facet has at least two nontriangular faces. If it
has exactly two, it is an antiprism.
\end{lemma}
\begin{proof}
If its only nontriangle has size $k$, then $d=k-3$, whereas the
antiprism bound gives $d\ge2(k-3)$, a contradiction. If the two sizes
are $k\ge l\ge4$, then
\[
 d=(k-3)+(l-3)\ge2(k-3).
\]
Thus $k=l$ and equality holds in the antiprism bound; its equality
statement identifies the facet with $A_k$.
\end{proof}

\begin{proof}[Proof of the lower bound in Theorem~\ref{thm:gap}]
Choose a facet $F$ of smallest positive defect $m\ge2$.
Lemma~\ref{lem:injection} gives $2W\ge m(1+q+s)$.
The following cases are exhaustive.
\begin{enumerate}
\item If $F=A_k$, then $m=2k-6$, $q=2$, and
$s=2k=m+6$ by Lemma~\ref{lem:antiprism}. Hence
$2W\ge m(m+9)\ge22$.
\item A non-antiprism with $m=2$ is impossible: $v_F=8$ gives
$k\le4$ for every face, so its defect is two quadrilaterals and
Lemma~\ref{lem:two} applies.
\item If $m=3$, Lemma~\ref{lem:nine} gives $2W\ge3(1+3+3)=21$.
\item If $m=4$, every face has size at most five. A pentagon gives
equality in the antiprism bound and has already been treated. Otherwise
there are four quadrilaterals, so $2W\ge4(1+4)=20$.
\item If $m\ge5$ and $F$ is not an antiprism,
Lemma~\ref{lem:two} gives $q\ge3$, whence $2W\ge4m\ge20$.
\end{enumerate}
Thus $W\ge10$. Since $W$ is divisible by three, $W\ge12$.
Equation~\eqref{eq:discrete} proves the lower bound. In particular the
three prospective intermediate defects $3,6,9$, with respective
volumes $5\pi^2/3$, $2\pi^2$ and $7\pi^2/3$, cannot occur.
\end{proof}

\section{Exact attainment by the double}
In Lorentz form $(-1,1,1,1,1)$, take the 24 unit normals $(1,a)$,
where $a$ runs through the $D_4$ roots having exactly two nonzero
entries, each $\pm1$. In the time-one Klein section their ideal
vertices are the eight $\pm e_i$ and the sixteen
$(\pm\tfrac12,\pm\tfrac12,\pm\tfrac12,\pm\tfrac12)$.
The inequalities are $|x_i|+|x_j|\le1$; they are bounded and their
origin is strict. Exact enumeration of all $\binom{24}{4}=10,626$
active four-wall subsets verifies that the displayed vertices are all
vertices of the intersection.

Reflect across $n=(1,1,1,0,0)$ by
$x\mapsto x-2\ip{x}{n}n$. Remove the gluing wall from both cells
and identify coincident walls. The resulting 38 halfspaces contain
the union. Conversely, on each side of the gluing hyperplane, their
inequalities together with that side's gluing inequality include all
inequalities of the corresponding original cell. Thus the new
intersection equals the union. In particular no unexamined escape
directions or additional vertices are being discarded.

The interiors are disjoint, so the volume is twice $4\pi^2/3$.
The exact face reconstruction supplies the following independent
incidence check.
\begin{center}
\begin{tabular}{lrr}
\toprule Quantity&24-cell&Double\\\midrule
Vertices&24&42\\ Edges&96&168\\ Ridges&96&164\\ Facets&24&38\\
Triangular ridges&96&152\\ Quadrilateral ridges&0&12\\
Octahedral facets&24&30\\ Doubled-octahedron facets&0&8\\
Ridge defect $W$&0&12\\ Reflection-orbifold Euler characteristic&1&2\\
\bottomrule
\end{tabular}
\end{center}
Each of the eight doubled-octahedron facets has nine vertices, eighteen
edges, eight triangles and three quadrilaterals. There are six shared
cusps and 36 other cusps. Every cusp has six incident facets, eight
edges and twelve ridges. The facet--vertex check is
$30\cdot6+8\cdot9=6\cdot42$. This completes attainment and
Theorem~\ref{thm:gap}.

\section{Verification and formalization scope}
The immutable written proof is
\path{results/ideal_ra4_second_volume_20261018.md}.
The producer
\path{scripts/fresh_ideal_ra4_double24_volume_20261018.py}
uses exact rational arithmetic to reconstruct the original and doubled
face lattices, verify Gram conditions at every ridge and cusp, and
compute all doubled-facet cusp ratios. The separate checker
\path{scripts/check_ideal_ra4_second_volume_saved_20261018.py}
imports no producer; it repeats the construction checks with a separate
elimination path and verifies the seven rational-function identities
used in the antiprism calculation. Generation, fresh replay, and the
final hardened checker all passed. Two independent agent reviewers
read the full written proof and both programs and checked the precise
published classification premises. They did not independently execute
the computation or reconstruct the historical classification census.

Neither program proves the universal geometric injection merely by
checking the double. In particular the construction JSON deliberately
does not identify its own computation with a proof of the second
minimum. The universal quantifiers are supplied by the written proof.

The Lean module \path{formal/CoxeterProofs/IdealRA4Gap.lean} contains
the finite-list budget and integer case-analysis interfaces: given
the geometric partition, facet case, and divisibility, the defect gap
and volume-numerator gap follow. Its geometric hypotheses are not
silently declared published theorems. The Lorentz/cusp injection,
shape identification and construction geometry are not yet an
end-to-end Lean development. The exact external/internal map and
reproduction commands are in
\path{docs/ideal_ra4_second_volume_lemma_map.md} and
\path{docs/ideal_ra4_second_volume_reproduce.md}.

\begin{thebibliography}{9}
\bibitem{Kolpakov} A. Kolpakov, \emph{On the optimality of the ideal
right-angled 24-cell}, Algebraic \& Geometric Topology \textbf{12}
(2012), 1941--1960. Published 27 October 2012.
\href{https://doi.org/10.2140/agt.2012.12.1941}{doi:10.2140/agt.2012.12.1941}.
\href{https://msp.org/agt/2012/12-4/agt-v12-n4-p03-s.pdf}{Published PDF}.
\end{thebibliography}
\end{document}
